\documentclass{article}
\usepackage{enumitem}
\usepackage{amsmath}
\begin{document}

\title{Chapter 18: Body Fluids and Circulation}
\date{}
\maketitle

\begin{enumerate}[label=Q\arabic*.]

\item The ABO blood group system is determined by:
\\(A) Presence or absence of Rh antigen on RBCs
\\(B) Presence of antigens (A and/or B) on RBC surface and corresponding antibodies in plasma
\\(C) Protein composition of plasma
\\(D) WBC surface markers

\item A person with blood group AB is called ``universal recipient'' because:
\\(A) They have both A and B antibodies in plasma
\\(B) Their plasma has neither anti-A nor anti-B antibodies, so they can receive blood from any ABO group
\\(C) They have no antigens on RBCs
\\(D) Their immune system destroys foreign RBCs efficiently

\item The Frank-Starling law of the heart states that:
\\(A) Heart rate increases with sympathetic stimulation
\\(B) Stroke volume increases with increased ventricular filling (preload) up to a physiological limit
\\(C) Cardiac output is fixed regardless of venous return
\\(D) Blood pressure is directly proportional to heart rate only

\item The sinoatrial (SA) node is called the pacemaker because:
\\(A) It generates nerve impulses from the brain
\\(B) It spontaneously generates electrical impulses that initiate each heartbeat at regular intervals
\\(C) It controls heart rate through hormonal signals
\\(D) It is located at the center of the heart

\item Assertion: Diastolic pressure represents the minimum blood pressure in arteries.\\
Reason: Diastolic pressure occurs during ventricular relaxation when the heart is refilling.
\\(A) Both true, Reason is correct explanation
\\(B) Both true, Reason is not correct explanation
\\(C) Assertion true, Reason false
\\(D) Assertion false, Reason true

\item The P wave on an ECG represents:
\\(A) Ventricular depolarization
\\(B) Atrial depolarization
\\(C) Ventricular repolarization
\\(D) AV node delay

\item Which of the following is correct about the QRS complex in ECG?
\\(A) Represents atrial repolarization
\\(B) Represents ventricular depolarization (spreading of impulse through ventricles)
\\(C) Represents atrial depolarization
\\(D) Represents ventricular repolarization

\item The lymphatic system returns lymph to blood circulation at the:
\\(A) Thoracic duct emptying into the left subclavian vein
\\(B) Portal vein
\\(C) Inferior vena cava
\\(D) Pulmonary artery

\item In a healthy adult, normal cardiac output at rest is approximately:
\\(A) 2--3 L/min
\\(B) 5--6 L/min
\\(C) 10--12 L/min
\\(D) 1--2 L/min

\item Which plasma protein is most important for maintaining oncotic (colloid osmotic) pressure of blood?
\\(A) Globulin
\\(B) Fibrinogen
\\(C) Albumin
\\(D) Transferrin

\item The AV node introduces a delay in impulse conduction to:
\\(A) Prevent atria from contracting simultaneously with ventricles
\\(B) Allow complete atrial contraction before ventricular contraction begins
\\(C) Increase heart rate
\\(D) Reduce blood pressure

\item Erythroblastosis fetalis (hemolytic disease of newborn) occurs when:
\\(A) Rh$^-$ mother carries Rh$^+$ fetus; maternal anti-Rh antibodies cross placenta and destroy fetal RBCs
\\(B) Rh$^+$ mother carries Rh$^-$ fetus
\\(C) ABO incompatibility between mother and fetus
\\(D) Mother lacks platelets

\item Blood coagulation requires which of the following in correct sequence?
\\(A) Platelet aggregation $\rightarrow$ Prothrombin $\rightarrow$ Thrombin $\rightarrow$ Fibrinogen $\rightarrow$ Fibrin
\\(B) Fibrinogen $\rightarrow$ Thrombin $\rightarrow$ Fibrin $\rightarrow$ Clot
\\(C) Thrombin $\rightarrow$ Prothrombin $\rightarrow$ Fibrin $\rightarrow$ Fibrinogen
\\(D) Platelet aggregation $\rightarrow$ Fibrin $\rightarrow$ Thrombin $\rightarrow$ Prothrombin

\item Match the following blood cells with their functions:\\
1. Neutrophils $\rightarrow$ (a) Phagocytosis of bacteria (first responders)\\
2. Basophils $\rightarrow$ (b) Release histamine in allergic reactions\\
3. Eosinophils $\rightarrow$ (c) Combat parasitic infections\\
4. Monocytes $\rightarrow$ (d) Differentiate into macrophages
\\(A) 1-a, 2-b, 3-c, 4-d
\\(B) 1-b, 2-a, 3-d, 4-c
\\(C) 1-c, 2-d, 3-a, 4-b
\\(D) 1-d, 2-c, 3-b, 4-a

\item Double circulation in humans means:
\\(A) Blood passes through the heart twice per heartbeat
\\(B) Blood circulates through two separate circuits: pulmonary (to lungs) and systemic (to body)
\\(C) Blood passes through two atria before reaching ventricles
\\(D) Both oxygenated and deoxygenated blood mix in the heart

\end{enumerate}
\end{document}
